\sectionbreak \section*{
  \gostTitleFont
  \redline
  ВВЕДЕНИЕ
}

\titlespace

\redline\gostFont
Современная розничная торговля предъявляет высокие требования к качеству
выкладки товаров на полках магазинов.
Соблюдение планограмм {--} утверждённых схем расположения продукции {--} напрямую
влияет на объём продаж, восприятие бренда покупателем и выполнение контрактных
обязательств между производителем и торговой сетью.
Компания «Санта Импэкс», являясь дистрибьютором широкого ассортимента
потребительских товаров, ежедневно направляет мерчендайзеров в сотни торговых
точек для контроля выкладки.
При этом процесс проверки по-прежнему остаётся преимущественно ручным:
сотрудник визуально оценивает полку, делает фотографию и вручную заполняет
отчёт, что влечёт за собой высокую трудоёмкость, субъективность оценки и
значительные временны́е затраты.

\redline\gostFont
Актуальность настоящей работы обусловлена необходимостью автоматизации
контроля мерчендайзинга с применением технологий компьютерного зрения
и глубокого обучения.
Развитие нейросетевых архитектур обнаружения объектов, в частности семейства
YOLO (You Only Look Once), позволяет в режиме реального времени анализировать
фотографии торговых полок, идентифицировать товарные позиции и сопоставлять
их расположение с эталонной планограммой.
Интеграция подобной системы с мессенджером Telegram даёт возможность
мерчендайзеру использовать привычный инструмент без установки дополнительного
программного обеспечения на рабочее устройство.

\redline\gostFont
Целью дипломного проекта является разработка программной системы
автоматизированной проверки выкладки товаров, включающей Telegram-бот
для полевых сотрудников и нейросетевой модуль на основе архитектуры YOLOv8,
обеспечивающий распознавание продукции и оценку соответствия планограмме.

\redline\gostFont
Для достижения поставленной цели необходимо решить следующие задачи:

\begin{itemize}[leftmargin=2.15cm, labelwidth=0.65cm, labelsep=0.0cm]
    \item[•] изучить предметную область мерчендайзинга и требования компании
          «Санта Импэкс» к контролю выкладки;
    \item[•] провести обзор существующих аналогов систем компьютерного зрения
          для ритейла;
    \item[•] спроектировать архитектуру системы: Telegram-бот, нейросетевой
          модуль обнаружения объектов и базу данных хранения результатов;
    \item[•] обучить модель YOLOv8 на размеченном наборе данных фотографий
          торговых полок с продукцией компании;
    \item[•] реализовать алгоритм сопоставления результатов детекции с
          планограммой и генерации аннотированного изображения;
    \item[•] разработать механизм автоматического формирования отчётов о
          корректности выкладки и их сохранения в базе данных;
    \item[•] провести тестирование системы и оценить метрики качества
          работы нейронной сети;
    \item[•] выполнить технико-экономическое обоснование разработки.
\end{itemize}

\redline\gostFont
Практическая значимость работы для компании «Санта Импэкс» определяется тремя ключевыми эффектами.
Во-первых, автоматизация значительно сокращает время оформления визита: мерчендайзер получает мгновенное заключение о корректности выкладки сразу после отправки фото в бот.
Во-вторых, централизованный сбор отчётов создаёт базу для системного анализа нарушений, что, наряду со снижением доли ручного труда, минимизирует операционные затраты и повышает эффективность работы персонала.
